% Modelo de Projeto de TCC, adaptado de "Modelo Canônico de Trabalho Acadêmico com abnTEX2".

\documentclass[
	% -- opções da classe memoir --
	12pt,				% tamanho da fonte
	openright,			% capítulos começam em pág ímpar (insere página vazia caso preciso)
	twoside,			% para impressão em verso e anverso. Oposto a oneside
	a4paper,			% tamanho do papel. 
	% -- opções da classe abntex2 --
	%chapter=TITLE,		% títulos de capítulos convertidos em letras maiúsculas
	%section=TITLE,		% títulos de seções convertidos em letras maiúsculas
	%subsection=TITLE,	% títulos de subseções convertidos em letras maiúsculas
	%subsubsection=TITLE,% títulos de subsubseções convertidos em letras maiúsculas
	% -- opções do pacote babel --
	english,			% idioma adicional para hifenização
	french,				% idioma adicional para hifenização
	spanish,			% idioma adicional para hifenização
	brazil,				% o último idioma é o principal do documento
	]{abntex2}

\input{config/pacotes_basicos.tex}

% ---
% Informações de dados para CAPA e FOLHA DE ROSTO
% ---
\titulo{Título do Projeto}
\autor{Nome do Autor}
\local{Brasília}
\data{3 de maio de 2019}
\orientador{Nome do orientador}
\coorientador{Nome do co-orientador}
\instituicao{%
  Universidade de Brasília -- UnB
  \par
  Instituto de Física}
\tipotrabalho{Monografia}
% O preambulo deve conter o tipo do trabalho, o objetivo, 
% o nome da instituição e a área de concentração 
\preambulo{Monografia apresentada ao Instituto de Física da Universidade de Brasília como parte dos requisitos  necessários à obtenção do título de Bacharel em Físicaısica.}
% ---

\input{config/layout.tex}
\input{config/ciencia.tex}

% ----
% Início do documento
% ----
\begin{document}

% Retira espaço extra obsoleto entre as frases.
\frenchspacing 

% ----------------------------------------------------------
% ELEMENTOS PRÉ-TEXTUAIS
% ----------------------------------------------------------
% \pretextual

% ---
% Capa
% ---
\imprimircapa
% ---

% ---
% Folha de rosto
% (o * indica que haverá a ficha bibliográfica)
% ---
\imprimirfolhaderosto*
% ---

% ---
% Inserir a ficha bibliografica
% ---
\input{texto/fichacatalografica.tex}
% ---

% ---
% Inserir errata
% ---
\input{texto/errata.tex}
% ---

% ---
% Inserir folha de aprovação
% ---
\input{texto/folhadeaprovacao.tex}
% ---

% ---
% Dedicatória
% ---
\input{texto/dedicatoria.tex}
% ---

% ---
% Agradecimentos
% ---
\input{texto/agradecimentos.tex}
% ---

% ---
% Epígrafe
% ---
\input{texto/epigrafe.tex}
% ---

% ---
% RESUMOS
% ---

% resumo em português
\input{texto/resumo.tex}

% resumo em inglês
\input{texto/abstract.tex}

% resumo em francês 
\input{texto/resume.tex}

% resumo em espanhol
\input{texto/resumen.tex}
% ---

% ---
% inserir lista de ilustrações
% ---
\pdfbookmark[0]{\listfigurename}{lof}
\listoffigures*
\cleardoublepage
% ---

% ---
% inserir lista de tabelas
% ---
\pdfbookmark[0]{\listtablename}{lot}
\listoftables*
\cleardoublepage
% ---

% ---
% inserir lista de abreviaturas e siglas
% ---
\input{texto/siglas.tex}
% ---

% ---
% inserir lista de símbolos
% ---
\input{texto/simbolos.tex}
% ---

% ---
% inserir o sumario
% ---
\pdfbookmark[0]{\contentsname}{toc}
\tableofcontents*
\cleardoublepage
% ---



% ----------------------------------------------------------
% ELEMENTOS TEXTUAIS
% ----------------------------------------------------------
\textual

% ----------------------------------------------------------
% Introdução
% ----------------------------------------------------------
\input{texto/introducao.tex}

% ----------------------------------------------------------
% PARTE - preparação da pesquisa
% ----------------------------------------------------------
\part{Título da parte I}

% ---
% Capitulos:
% ---

\input{texto/capitulo-1.tex}

\input{texto/capitulo-2.tex}

% ----------------------------------------------------------
% Parte de revisão de literatura
% ----------------------------------------------------------
\part{Título da parte II}

% ---
% Capitulos:
% ---

\input{texto/capitulo-3.tex}

\input{texto/capitulo-4.tex}


% ----------------------------------------------------------
% Resultados
% ----------------------------------------------------------
\part{Título da parte III}

% ---
% Capitulos:
% ---

\input{texto/capitulo-5.tex}

\input{texto/capitulo-6.tex}

% ---
% Finaliza a parte no bookmark do PDF, para que se inicie o bookmark na raiz
% ---
\bookmarksetup{startatroot}% 
% ---

% ---
% Conclusão
% ---
% ---
% Capitulos:
% ---

\input{texto/conclusao.tex}

% ----------------------------------------------------------
% ELEMENTOS PÓS-TEXTUAIS
% ----------------------------------------------------------
\postextual


% ----------------------------------------------------------
% Referências bibliográficas
% ----------------------------------------------------------
\bibliography{monografia.bib}

% ----------------------------------------------------------
% Glossário
% ----------------------------------------------------------
%
% Consulte o manual da classe abntex2 para orientações sobre o glossário.
%
%\glossary

% ----------------------------------------------------------
% Apêndices
% ----------------------------------------------------------

% ---
% Inicia os apêndices
% ---
\begin{apendicesenv}

% Imprime uma página indicando o início dos apêndices
\partapendices

\input{texto/apendice-1.tex}

\input{texto/apendice-2.tex}


\end{apendicesenv}
% ---


% ----------------------------------------------------------
% Anexos
% ----------------------------------------------------------

% ---
% Inicia os anexos
% ---
\begin{anexosenv}

% Imprime uma página indicando o início dos anexos
\partanexos

\input{texto/anexo-1.tex}

\input{texto/anexo-2.tex}

\end{anexosenv}

%---------------------------------------------------------------------
% INDICE REMISSIVO
%---------------------------------------------------------------------

\printindex

\end{document}

